\section{生物・遺伝}

\subsection{平成27年度 問16}
\label{q:h27-reg-applied-16}

\noindent\textbf{原典：} sources/H27/17\_applied\_science.pdf、PDF物理6頁。

\textbf{公式正答：} 2（single）。

\par\textit{原典本文の校正と独自解説は未作成。分類のみの登録項目。}

\subsection{平成30年度 問15}
\label{q:h30-reg-applied-15}

\noindent\textbf{原典：} sources/H30/17\_applied\_science.pdf、PDF物理5頁。

\textbf{公式正答：} 5（single）。

\par\textit{原典本文の校正と独自解説は未作成。分類のみの登録項目。}

\section{有機・高分子}

\subsection{令和7年度 問11}
\label{q:r07-reg-applied-11}

\noindent\textbf{原典：} sources/R07/17\_applied\_science.pdf、PDF物理4頁。

\textbf{公式正答：} 5（single）。

\par\textit{原典本文の校正と独自解説は未作成。分類のみの登録項目。}

\section{芳香族求電子置換の配向性と導入順}

\subsection{平成25年度 問17}
\label{q:h25-reg-applied-17}

\noindent\textbf{原典：} sources/H25/17\_applied\_science.pdf、PDF物理7頁。

\textbf{公式正答：} 4（single）。

\paragraph{問題の要旨} 芳香族置換の配向性からoニトロトルエンの導入順を決める。

\Needspace{4\baselineskip}\paragraph{確認した条件・記述（独自要約）}\mbox{}\par

\noindent\textbf{判断対象：} 芳香族置換の配向性からoニトロトルエンの導入順を決める。\par

\noindent\textbf{原典確認：} 化学式・否定・表の全5候補を保存原典24頁画像で照合。\par

\paragraph{問われる核心} 芳香族置換の配向性からoニトロトルエンの導入順を決める。

\paragraph{必要知識} メチル基は環へ電子を供与してo,p位の求電子置換を促し、ニトロ基は電子を吸引してm配向となる。

\paragraph{考え方・正答を選ぶ理由（査読待ち初稿）} メチル基は環へ電子を供与してo,p位の求電子置換を促し、ニトロ基は電子を吸引してm配向となる。先にメチルを入れてトルエンを作り、その後ニトロを導入する。イだけをベンゼンへ入れたものはニトロベンゼン。この四対応で候補4。この経路でもo体だけが100\%得られるという意味ではなく、他位置との生成混合を分離する。

\Needspace{4\baselineskip}\paragraph{選択肢の検討（査読待ち初稿）}\mbox{}\par

\noindent\textbf{1：選ばない候補} ニトロ基/メチル基/ニトロベンゼン/トルエン。ニトロを先に入れる順はm配向・不活性化なので目的の経路に合わない。\par

\noindent\textbf{2：選ばない候補} カルボキシル基/ニトロ基/安息香酸/ニトロベンゼン。oニトロトルエンの置換基メチル/ニトロ又は中間トルエン/ニトロベンゼンの対応と異なる。\par

\noindent\textbf{3：選ばない候補} ニトロ基/アミノ基/ニトロベンゼン/アニリン。oニトロトルエンの置換基メチル/ニトロ又は中間トルエン/ニトロベンゼンの対応と異なる。\par

\noindent\textbf{4：選ぶ候補} メチル基/ニトロ基/トルエン/ニトロベンゼン。メチル基はo,p配向、トルエンを先に作りニトロを導入する全対応に合う。\par

\noindent\textbf{5：選ばない候補} ヒドロキシ基/ニトロ基/フェノール/ニトロベンゼン。oニトロトルエンの置換基メチル/ニトロ又は中間トルエン/ニトロベンゼンの対応と異なる。\par

\paragraph{再利用する考え方} 物質名を暗記する前に、結合の構造、反応の種類、成立条件、数量の保存を対応させる。

\paragraph{共通理論} \hyperref[theory:organic:directing]{第\ref*{theory:organic:directing}節}

\par 参考：\cite{organic_directing}。

\par 難易度 2（校準初稿）、優先度未評価。

\subsection{令和元年度再試験 問13}
\label{q:r01-ret-applied-13}

\noindent\textbf{原典：} sources/R01\_retest/17\_applied\_science.pdf、PDF物理5頁。

\textbf{公式正答：} 1（single）。

\paragraph{問題の要旨} 芳香族置換の配向性からoニトロトルエンの導入順を決める。

\Needspace{4\baselineskip}\paragraph{確認した条件・記述（独自要約）}\mbox{}\par

\noindent\textbf{判断対象：} 芳香族置換の配向性からoニトロトルエンの導入順を決める。\par

\noindent\textbf{原典確認：} 化学式・否定・表の全5候補を保存原典24頁画像で照合。\par

\paragraph{問われる核心} 芳香族置換の配向性からoニトロトルエンの導入順を決める。

\paragraph{必要知識} メチル基は環へ電子を供与してo,p位の求電子置換を促し、ニトロ基は電子を吸引してm配向となる。

\paragraph{考え方・正答を選ぶ理由（査読待ち初稿）} メチル基は環へ電子を供与してo,p位の求電子置換を促し、ニトロ基は電子を吸引してm配向となる。先にメチルを入れてトルエンを作り、その後ニトロを導入する。イだけをベンゼンへ入れたものはニトロベンゼン。この四対応で候補1。この経路でもo体だけが100\%得られるという意味ではなく、他位置との生成混合を分離する。

\Needspace{4\baselineskip}\paragraph{選択肢の検討（査読待ち初稿）}\mbox{}\par

\noindent\textbf{1：選ぶ候補} メチル基/ニトロ基/トルエン/ニトロベンゼン。メチル基はo,p配向、トルエンを先に作りニトロを導入する全対応に合う。\par

\noindent\textbf{2：選ばない候補} ニトロ基/カルボキシル基/ニトロベンゼン/安息香酸。oニトロトルエンの置換基メチル/ニトロ又は中間トルエン/ニトロベンゼンの対応と異なる。\par

\noindent\textbf{3：選ばない候補} アミノ基/ニトロ基/アニリン/ニトロベンゼン。oニトロトルエンの置換基メチル/ニトロ又は中間トルエン/ニトロベンゼンの対応と異なる。\par

\noindent\textbf{4：選ばない候補} ニトロ基/メチル基/ニトロベンゼン/トルエン。ニトロを先に入れる順はm配向・不活性化なので目的の経路に合わない。\par

\noindent\textbf{5：選ばない候補} ニトロ基/ヒドロキシ基/ニトロベンゼン/フェノール。oニトロトルエンの置換基メチル/ニトロ又は中間トルエン/ニトロベンゼンの対応と異なる。\par

\paragraph{再利用する考え方} 物質名を暗記する前に、結合の構造、反応の種類、成立条件、数量の保存を対応させる。

\paragraph{共通理論} \hyperref[theory:organic:directing]{第\ref*{theory:organic:directing}節}

\par 参考：\cite{organic_directing}。

\par 難易度 2（校準初稿）、優先度未評価。

\section{ベンゼン非局在πと単置換対称性}

\subsection{令和元年度 問11}
\label{q:r01-reg-applied-11}

\noindent\textbf{原典：} sources/R01/17\_applied\_science.pdf、PDF物理4頁。

\textbf{公式正答：} 3（single）。

\paragraph{問題の要旨} ベンゼンの非局在π電子と単置換位置の対称性を判定する。

\Needspace{4\baselineskip}\paragraph{確認した条件・記述（独自要約）}\mbox{}\par

\noindent\textbf{判断対象：} ベンゼンの非局在π電子と単置換位置の対称性を判定する。\par

\noindent\textbf{原典確認：} 化学式・否定・表の全5候補を保存原典24頁画像で照合。\par

\paragraph{問われる核心} ベンゼンの非局在π電子と単置換位置の対称性を判定する。

\paragraph{必要知識} 六炭素のp軌道が環全体で重なり、6π電子が非局在して等価なC-C結合を作る。

\paragraph{考え方・正答を選ぶ理由（査読待ち初稿）} 六炭素のp軌道が環全体で重なり、6π電子が非局在して等価なC-C結合を作る。ベンゼンの六つのH位置は等価なので一つのBrを置いたモノブロモベンゼンに位置異性体はなく3。二置換では相対位置o,m,pを区別する。共鳴式を実際に単/二重結合が交互に動く混合物と考えない。

\Needspace{4\baselineskip}\paragraph{選択肢の検討（査読待ち初稿）}\mbox{}\par

\noindent\textbf{1：選ばない候補} C6H6の六炭素環。各炭素がHを一つ持つ平面環。\par

\noindent\textbf{2：選ばない候補} 通常アルケンより安定・反応しにくい。芳香族性により通常の付加反応へ抵抗し、置換で環共役を回復する。全反応速度を無条件比較しない。\par

\noindent\textbf{3：選ぶ候補} モノブロモベンゼンに位置異性体。六位置が対称で等価なので単置換の位置異性体は一種のみ。\par

\noindent\textbf{4：選ばない候補} 全C-C結合は単/二重の中間。非局在π電子により等価な結合となる。\par

\noindent\textbf{5：選ばない候補} C-C結合にπ電子が関与。各sp2炭素のp軌道から環全体へ6π電子を持つ。\par

\paragraph{再利用する考え方} 物質名を暗記する前に、結合の構造、反応の種類、成立条件、数量の保存を対応させる。

\paragraph{共通理論} \hyperref[theory:organic:benzene]{第\ref*{theory:organic:benzene}節}

\par 参考：\cite{organic_benzene}。

\par 難易度 2（校準初稿）、優先度未評価。

\subsection{令和4年度 問18}
\label{q:r04-reg-applied-18}

\noindent\textbf{原典：} sources/R04/17\_applied\_science.pdf、PDF物理6頁。

\textbf{公式正答：} 4（single）。

\paragraph{問題の要旨} ベンゼンの非局在π電子と単置換位置の対称性を判定する。

\Needspace{4\baselineskip}\paragraph{確認した条件・記述（独自要約）}\mbox{}\par

\noindent\textbf{判断対象：} ベンゼンの非局在π電子と単置換位置の対称性を判定する。\par

\noindent\textbf{原典確認：} 化学式・否定・表の全5候補を保存原典24頁画像で照合。\par

\paragraph{問われる核心} ベンゼンの非局在π電子と単置換位置の対称性を判定する。

\paragraph{必要知識} 六炭素のp軌道が環全体で重なり、6π電子が非局在して等価なC-C結合を作る。

\paragraph{考え方・正答を選ぶ理由（査読待ち初稿）} 六炭素のp軌道が環全体で重なり、6π電子が非局在して等価なC-C結合を作る。ベンゼンの六つのH位置は等価なので一つのBrを置いたモノブロモベンゼンに位置異性体はなく4。二置換では相対位置o,m,pを区別する。共鳴式を実際に単/二重結合が交互に動く混合物と考えない。

\Needspace{4\baselineskip}\paragraph{選択肢の検討（査読待ち初稿）}\mbox{}\par

\noindent\textbf{1：選ばない候補} 通常アルケンより安定・反応しにくい。芳香族性により通常の付加反応へ抵抗し、置換で環共役を回復する。全反応速度を無条件比較しない。\par

\noindent\textbf{2：選ばない候補} C-C結合にπ電子が関与。各sp2炭素のp軌道から環全体へ6π電子を持つ。\par

\noindent\textbf{3：選ばない候補} 全C-C結合は単/二重の中間。非局在π電子により等価な結合となる。\par

\noindent\textbf{4：選ぶ候補} モノブロモベンゼンに位置異性体。六位置が対称で等価なので単置換の位置異性体は一種のみ。\par

\noindent\textbf{5：選ばない候補} C6H6の六炭素環。各炭素がHを一つ持つ平面環。\par

\paragraph{再利用する考え方} 物質名を暗記する前に、結合の構造、反応の種類、成立条件、数量の保存を対応させる。

\paragraph{共通理論} \hyperref[theory:organic:benzene]{第\ref*{theory:organic:benzene}節}

\par 参考：\cite{organic_benzene}。

\par 難易度 2（校準初稿）、優先度未評価。

\section{主鎖・環境条件と生分解性}

\subsection{平成24年度 問15}
\label{q:h24-reg-applied-15}

\noindent\textbf{原典：} sources/H24/17\_applied\_science.pdf、PDF物理5頁。

\textbf{公式正答：} 4（single）。

\paragraph{問題の要旨} 主鎖結合と生分解性高分子の代表分類を判定する。

\Needspace{4\baselineskip}\paragraph{確認した条件・記述（独自要約）}\mbox{}\par

\noindent\textbf{判断対象：} 主鎖結合と生分解性高分子の代表分類を判定する。\par

\noindent\textbf{原典確認：} 化学式・否定・表の全5候補を保存原典24頁画像で照合。\par

\paragraph{問われる核心} 主鎖結合と生分解性高分子の代表分類を判定する。

\paragraph{必要知識} PGA、PLA、PBS、PHBは代表的な生分解性ポリエステルである。

\paragraph{考え方・正答を選ぶ理由（査読待ち初稿）} PGA、PLA、PBS、PHBは代表的な生分解性ポリエステルである。PMMAは側鎖にエステルを持つが主鎖はC-Cで、同じ分類に入らないので4。エステルという単語だけで生分解を保証せず、主鎖・分解環境・試験条件を区別する。生分解性は海洋でも直ちに消えるという意味ではない。

\Needspace{4\baselineskip}\paragraph{選択肢の検討（査読待ち初稿）}\mbox{}\par

\noindent\textbf{1：選ばない候補} PGA。脂肪族ポリエステルとして生分解性分類。\par

\noindent\textbf{2：選ばない候補} ポリL乳酸。PLAの生分解性で分類される。\par

\noindent\textbf{3：選ばない候補} PBS。PBSは脂肪族ポリエステル。\par

\noindent\textbf{4：選ぶ候補} PMMA。炭素主鎖の通常PMMAは生分解性高分子に該当しない。\par

\noindent\textbf{5：選ばない候補} ポリR3ヒドロキシブチレート。PHBは微生物が作るポリエステル。\par

\paragraph{再利用する考え方} 物質名を暗記する前に、結合の構造、反応の種類、成立条件、数量の保存を対応させる。

\paragraph{共通理論} \hyperref[theory:organic:biodegradation]{第\ref*{theory:organic:biodegradation}節}

\par 参考：\cite{organic_jbpa,organic_jbpa_en}。

\par 難易度 2（校準初稿）、優先度未評価。

\section{炭素骨格・脂肪酸と構造異性体}

\subsection{平成23年度 問15}
\label{q:h23-reg-applied-15}

\noindent\textbf{原典：} sources/H23/17\_applied\_science.pdf、PDF物理6頁。

\textbf{公式正答：} 2（single）。

\paragraph{問題の要旨} 18炭素の飽和脂肪酸を縮合構造式で数える。

\Needspace{4\baselineskip}\paragraph{確認した条件・記述（独自要約）}\mbox{}\par

\noindent\textbf{判断対象：} 18炭素の飽和脂肪酸を縮合構造式で数える。\par

\noindent\textbf{原典確認：} 化学式・否定・表の全5候補を保存原典24頁画像で照合。\par

\paragraph{問われる核心} 18炭素の飽和脂肪酸を縮合構造式で数える。

\paragraph{必要知識} オクタデカン酸は18炭素の飽和一価カルボン酸である。

\paragraph{考え方・正答を選ぶ理由（査読待ち初稿）} オクタデカン酸は18炭素の飽和一価カルボン酸である。CH3の1、(CH2)の個数、COOHの1を合わせるので中間は16、2。4と5は炭素数18でも不飽和でステアリン酸ではない。一般式C18H36O2も確認できる。

\Needspace{4\baselineskip}\paragraph{選択肢の検討（査読待ち初稿）}\mbox{}\par

\noindent\textbf{1：選ばない候補} CH3(CH2)14COOH。16炭素の飽和酸。\par

\noindent\textbf{2：選ぶ候補} CH3(CH2)16COOH。1+16+1=\allowbreak{}18炭素で飽和。\par

\noindent\textbf{3：選ばない候補} CH3(CH2)18COOH。20炭素の飽和酸。\par

\noindent\textbf{4：選ばない候補} CH3(CH2)7CH=\allowbreak{}CH(CH2)7COOH。18炭素だが二重結合1個。\par

\noindent\textbf{5：選ばない候補} CH3(CH2)4CH=\allowbreak{}CHCH2CH=\allowbreak{}CH(CH2)7COOH。18炭素だが二重結合2個。\par

\paragraph{再利用する考え方} 物質名を暗記する前に、結合の構造、反応の種類、成立条件、数量の保存を対応させる。

\paragraph{共通理論} \hyperref[theory:organic:structure]{第\ref*{theory:organic:structure}節}

\par 難易度 3（校準初稿）、優先度未評価。

\subsection{令和元年度再試験 問11}
\label{q:r01-ret-applied-11}

\noindent\textbf{原典：} sources/R01\_retest/17\_applied\_science.pdf、PDF物理4頁。

\textbf{公式正答：} 4（single）。

\paragraph{問題の要旨} C4からC6アルカンの骨格異性体を重複せず列挙する。

\Needspace{4\baselineskip}\paragraph{確認した条件・記述（独自要約）}\mbox{}\par

\noindent\textbf{判断対象：} C4からC6アルカンの骨格異性体を重複せず列挙する。\par

\noindent\textbf{原典確認：} 化学式・否定・表の全5候補を保存原典24頁画像で照合。\par

\paragraph{問われる核心} C4からC6アルカンの骨格異性体を重複せず列挙する。

\paragraph{必要知識} 結合のつながりが違う構造異性体を数える。

\paragraph{考え方・正答を選ぶ理由（査読待ち初稿）} 結合のつながりが違う構造異性体を数える。C4はnブタンと2メチルプロパンの2。C5はnペンタン、2メチルブタン、2,2ジメチルプロパンの3。C6はnヘキサン、2/3メチルペンタン、2,2/2,3ジメチルブタンの5。ゆえに4。鎖を逆から番号付けした同じ骨格を二回数えない。

\Needspace{4\baselineskip}\paragraph{選択肢の検討（査読待ち初稿）}\mbox{}\par

\noindent\textbf{1：選ばない候補} 1/2/3。三つとも骨格数を少なく数えている。\par

\noindent\textbf{2：選ばない候補} 1/3/4。C4とC6を少なく数える。\par

\noindent\textbf{3：選ばない候補} 2/3/4。C6は5であって4ではない。\par

\noindent\textbf{4：選ぶ候補} 2/3/5。C4=\allowbreak{}2,C5=\allowbreak{}3,C6=\allowbreak{}5の全対応。\par

\noindent\textbf{5：選ばない候補} 3/4/5。C4とC5を一つずつ多く数える。\par

\paragraph{再利用する考え方} 物質名を暗記する前に、結合の構造、反応の種類、成立条件、数量の保存を対応させる。

\paragraph{共通理論} \hyperref[theory:organic:structure]{第\ref*{theory:organic:structure}節}

\par 参考：\cite{organic_isomers}。

\par 難易度 2（校準初稿）、優先度未評価。

\section{燃焼元素分析と候補分子式}

\subsection{平成28年度 問15}
\label{q:h28-reg-applied-15}

\noindent\textbf{原典：} sources/H28/17\_applied\_science.pdf、PDF物理6頁。

\textbf{公式正答：} 5（single）。

\par\textbf{検討未確定：} 原典は酸素を含むのに芳香族炭化水素と呼ぶ。炭化水素はC,Hのみで名称に疑義がある。数値と候補は保持し、C,H,O有機化合物として条件付きで解く。

\paragraph{問題の要旨} 燃焼生成物から元素量を保存して有機化合物を選ぶ。

\Needspace{4\baselineskip}\paragraph{確認した条件・記述（独自要約）}\mbox{}\par

\noindent\textbf{試料：} 69mg\par

\noindent\textbf{CO2：} 154mg\par

\noindent\textbf{H2O：} 27mg\par

\noindent\textbf{原子量：} C=\allowbreak{}12,H=\allowbreak{}1,O=\allowbreak{}16\par

\paragraph{問われる核心} 燃焼生成物から元素量を保存して有機化合物を選ぶ。

\paragraph{必要知識} 完全燃焼で試料CはCO2、試料HはH2Oへ全て移る。

\paragraph{考え方・正答を選ぶ理由（査読待ち初稿）} 完全燃焼で試料CはCO2、試料HはH2Oへ全て移る。炭素質量154×12/44=\allowbreak{}42.0mg、水素質量27×2/18=\allowbreak{}3.0mg。試料がC,H,Oのみなので酸素質量は差の24.0mg。各質量を12,1,16で割った原子比は7:6:3となる。候補中C7H6O3のサリチル酸だけが合い5。組成比だけでは一般の分子式を一意にしないが、提示候補比較で選べる。

\Needspace{4\baselineskip}\paragraph{選択肢の検討（査読待ち初稿）}\mbox{}\par

\noindent\textbf{1：選ばない候補} トルエン。C7H8の原子数比が導出比に一致しない。\par

\noindent\textbf{2：選ばない候補} フェノール。C6H6Oの原子数比が導出比に一致しない。\par

\noindent\textbf{3：選ばない候補} 安息香酸。C7H6O2の原子数比が導出比に一致しない。\par

\noindent\textbf{4：選ばない候補} テレフタル酸。C8H6O4の原子数比が導出比に一致しない。\par

\noindent\textbf{5：選ぶ候補} サリチル酸。C7H6O3の原子数比が導出比に一致する。\par

\paragraph{再利用する考え方} 物質名を暗記する前に、結合の構造、反応の種類、成立条件、数量の保存を対応させる。

\paragraph{共通理論} \hyperref[theory:organic:combustion]{第\ref*{theory:organic:combustion}節}

\par 難易度 3（校準初稿）、優先度未評価。

\subsection{平成29年度 問15}
\label{q:h29-reg-applied-15}

\noindent\textbf{原典：} sources/H29/17\_applied\_science.pdf、PDF物理5頁。

\textbf{公式正答：} 2（single）。

\paragraph{問題の要旨} 燃焼生成物から元素量を保存して有機化合物を選ぶ。

\Needspace{4\baselineskip}\paragraph{確認した条件・記述（独自要約）}\mbox{}\par

\noindent\textbf{試料：} 88mg\par

\noindent\textbf{CO2：} 176mg\par

\noindent\textbf{H2O：} 72mg\par

\noindent\textbf{原子量：} C=\allowbreak{}12,H=\allowbreak{}1,O=\allowbreak{}16\par

\paragraph{問われる核心} 燃焼生成物から元素量を保存して有機化合物を選ぶ。

\paragraph{必要知識} 完全燃焼で試料CはCO2、試料HはH2Oへ全て移る。

\paragraph{考え方・正答を選ぶ理由（査読待ち初稿）} 完全燃焼で試料CはCO2、試料HはH2Oへ全て移る。炭素質量176×12/44=\allowbreak{}48.0mg、水素質量72×2/18=\allowbreak{}8.0mg。試料がC,H,Oのみなので酸素質量は差の32.0mg。各質量を12,1,16で割った原子比は2:4:1となる。候補中C2H4Oのアセトアルデヒドだけが合い2。組成比だけでは一般の分子式を一意にしないが、提示候補比較で選べる。

\Needspace{4\baselineskip}\paragraph{選択肢の検討（査読待ち初稿）}\mbox{}\par

\noindent\textbf{1：選ばない候補} エタノール。C2H6Oの原子数比が導出比に一致しない。\par

\noindent\textbf{2：選ぶ候補} アセトアルデヒド。C2H4Oの原子数比が導出比に一致する。\par

\noindent\textbf{3：選ばない候補} アセトン。C3H6Oの原子数比が導出比に一致しない。\par

\noindent\textbf{4：選ばない候補} 酢酸。C2H4O2の原子数比が導出比に一致しない。\par

\noindent\textbf{5：選ばない候補} アクリル酸。C3H4O2の原子数比が導出比に一致しない。\par

\paragraph{再利用する考え方} 物質名を暗記する前に、結合の構造、反応の種類、成立条件、数量の保存を対応させる。

\paragraph{共通理論} \hyperref[theory:organic:combustion]{第\ref*{theory:organic:combustion}節}

\par 難易度 3（校準初稿）、優先度未評価。

\subsection{令和元年度 問13}
\label{q:r01-reg-applied-13}

\noindent\textbf{原典：} sources/R01/17\_applied\_science.pdf、PDF物理4頁。

\textbf{公式正答：} 3（single）。

\paragraph{問題の要旨} 燃焼生成物から元素量を保存して有機化合物を選ぶ。

\Needspace{4\baselineskip}\paragraph{確認した条件・記述（独自要約）}\mbox{}\par

\noindent\textbf{試料：} 61mg\par

\noindent\textbf{CO2：} 154mg\par

\noindent\textbf{H2O：} 27mg\par

\noindent\textbf{原子量：} C=\allowbreak{}12,H=\allowbreak{}1,O=\allowbreak{}16\par

\paragraph{問われる核心} 燃焼生成物から元素量を保存して有機化合物を選ぶ。

\paragraph{必要知識} 完全燃焼で試料CはCO2、試料HはH2Oへ全て移る。

\paragraph{考え方・正答を選ぶ理由（査読待ち初稿）} 完全燃焼で試料CはCO2、試料HはH2Oへ全て移る。炭素質量154×12/44=\allowbreak{}42.0mg、水素質量27×2/18=\allowbreak{}3.0mg。試料がC,H,Oのみなので酸素質量は差の16.0mg。各質量を12,1,16で割った原子比は7:6:2となる。候補中C7H6O2の安息香酸だけが合い3。組成比だけでは一般の分子式を一意にしないが、提示候補比較で選べる。

\Needspace{4\baselineskip}\paragraph{選択肢の検討（査読待ち初稿）}\mbox{}\par

\noindent\textbf{1：選ばない候補} トルエン。C7H8の原子数比が導出比に一致しない。\par

\noindent\textbf{2：選ばない候補} フェノール。C6H6Oの原子数比が導出比に一致しない。\par

\noindent\textbf{3：選ぶ候補} 安息香酸。C7H6O2の原子数比が導出比に一致する。\par

\noindent\textbf{4：選ばない候補} テレフタル酸。C8H6O4の原子数比が導出比に一致しない。\par

\noindent\textbf{5：選ばない候補} サリチル酸。C7H6O3の原子数比が導出比に一致しない。\par

\paragraph{再利用する考え方} 物質名を暗記する前に、結合の構造、反応の種類、成立条件、数量の保存を対応させる。

\paragraph{共通理論} \hyperref[theory:organic:combustion]{第\ref*{theory:organic:combustion}節}

\par 難易度 3（校準初稿）、優先度未評価。

\subsection{令和5年度 問11}
\label{q:r05-reg-applied-11}

\noindent\textbf{原典：} sources/R05/17\_applied\_science.pdf、PDF物理4頁。

\textbf{公式正答：} 4（single）。

\paragraph{問題の要旨} 燃焼生成物から元素量を保存して有機化合物を選ぶ。

\Needspace{4\baselineskip}\paragraph{確認した条件・記述（独自要約）}\mbox{}\par

\noindent\textbf{試料：} 92mg\par

\noindent\textbf{CO2：} 176mg\par

\noindent\textbf{H2O：} 108mg\par

\noindent\textbf{原子量：} C=\allowbreak{}12,H=\allowbreak{}1,O=\allowbreak{}16\par

\paragraph{問われる核心} 燃焼生成物から元素量を保存して有機化合物を選ぶ。

\paragraph{必要知識} 完全燃焼で試料CはCO2、試料HはH2Oへ全て移る。

\paragraph{考え方・正答を選ぶ理由（査読待ち初稿）} 完全燃焼で試料CはCO2、試料HはH2Oへ全て移る。炭素質量176×12/44=\allowbreak{}48.0mg、水素質量108×2/18=\allowbreak{}12.0mg。試料がC,H,Oのみなので酸素質量は差の32.0mg。各質量を12,1,16で割った原子比は2:6:1となる。候補中C2H6Oのエタノールだけが合い4。組成比だけでは一般の分子式を一意にしないが、提示候補比較で選べる。

\Needspace{4\baselineskip}\paragraph{選択肢の検討（査読待ち初稿）}\mbox{}\par

\noindent\textbf{1：選ばない候補} アクリル酸。C3H4O2の原子数比が導出比に一致しない。\par

\noindent\textbf{2：選ばない候補} アセトアルデヒド。C2H4Oの原子数比が導出比に一致しない。\par

\noindent\textbf{3：選ばない候補} アセトン。C3H6Oの原子数比が導出比に一致しない。\par

\noindent\textbf{4：選ぶ候補} エタノール。C2H6Oの原子数比が導出比に一致する。\par

\noindent\textbf{5：選ばない候補} 酢酸。C2H4O2の原子数比が導出比に一致しない。\par

\paragraph{再利用する考え方} 物質名を暗記する前に、結合の構造、反応の種類、成立条件、数量の保存を対応させる。

\paragraph{共通理論} \hyperref[theory:organic:combustion]{第\ref*{theory:organic:combustion}節}

\par 難易度 3（校準初稿）、優先度未評価。

\subsection{令和7年度 問12}
\label{q:r07-reg-applied-12}

\noindent\textbf{原典：} sources/R07/17\_applied\_science.pdf、PDF物理4頁。

\textbf{公式正答：} 2（single）。

\par\textbf{検討未確定：} 原典は酸素を含むのに芳香族炭化水素と呼ぶ。炭化水素はC,Hのみで名称に疑義がある。数値と候補は保持し、C,H,O有機化合物として条件付きで解く。

\paragraph{問題の要旨} 燃焼生成物から元素量を保存して有機化合物を選ぶ。

\Needspace{4\baselineskip}\paragraph{確認した条件・記述（独自要約）}\mbox{}\par

\noindent\textbf{試料：} 47mg\par

\noindent\textbf{CO2：} 132mg\par

\noindent\textbf{H2O：} 27mg\par

\noindent\textbf{原子量：} C=\allowbreak{}12,H=\allowbreak{}1,O=\allowbreak{}16\par

\paragraph{問われる核心} 燃焼生成物から元素量を保存して有機化合物を選ぶ。

\paragraph{必要知識} 完全燃焼で試料CはCO2、試料HはH2Oへ全て移る。

\paragraph{考え方・正答を選ぶ理由（査読待ち初稿）} 完全燃焼で試料CはCO2、試料HはH2Oへ全て移る。炭素質量132×12/44=\allowbreak{}36.0mg、水素質量27×2/18=\allowbreak{}3.0mg。試料がC,H,Oのみなので酸素質量は差の8.0mg。各質量を12,1,16で割った原子比は6:6:1となる。候補中C6H6Oのフェノールだけが合い2。組成比だけでは一般の分子式を一意にしないが、提示候補比較で選べる。

\Needspace{4\baselineskip}\paragraph{選択肢の検討（査読待ち初稿）}\mbox{}\par

\noindent\textbf{1：選ばない候補} トルエン。C7H8の原子数比が導出比に一致しない。\par

\noindent\textbf{2：選ぶ候補} フェノール。C6H6Oの原子数比が導出比に一致する。\par

\noindent\textbf{3：選ばない候補} 安息香酸。C7H6O2の原子数比が導出比に一致しない。\par

\noindent\textbf{4：選ばない候補} テレフタル酸。C8H6O4の原子数比が導出比に一致しない。\par

\noindent\textbf{5：選ばない候補} サリチル酸。C7H6O3の原子数比が導出比に一致しない。\par

\paragraph{再利用する考え方} 物質名を暗記する前に、結合の構造、反応の種類、成立条件、数量の保存を対応させる。

\paragraph{共通理論} \hyperref[theory:organic:combustion]{第\ref*{theory:organic:combustion}節}

\par 難易度 3（校準初稿）、優先度未評価。

\section{汎用樹脂・エンプラとグレード条件}

\subsection{令和元年度再試験 問17}
\label{q:r01-ret-applied-17}

\noindent\textbf{原典：} sources/R01\_retest/17\_applied\_science.pdf、PDF物理6頁。

\textbf{公式正答：} 1（single）。

\par\textbf{検討未確定：} 原典はPEのグレードを限定しない。汎用PEを想定した公式1を示し、UHMWPE等の工学材料例と区別。

\paragraph{問題の要旨} 汎用樹脂と代表エンジニアリングプラスチックを区別する。

\Needspace{4\baselineskip}\paragraph{確認した条件・記述（独自要約）}\mbox{}\par

\noindent\textbf{判断対象：} 汎用樹脂と代表エンジニアリングプラスチックを区別する。\par

\noindent\textbf{原典確認：} 化学式・否定・表の全5候補を保存原典24頁画像で照合。\par

\paragraph{問われる核心} 汎用樹脂と代表エンジニアリングプラスチックを区別する。

\paragraph{必要知識} エンプラは耐熱・機械特性等により部品へ使う実用分類である。

\paragraph{考え方・正答を選ぶ理由（査読待ち初稿）} エンプラは耐熱・機械特性等により部品へ使う実用分類である。PC、PA、POM、PPO/PPEは代表例で、通常PEは汎用樹脂として1を選ぶ。分類は全材料を厳密な一数値で区切るものではなく、超高分子量PE等は工学材料として使われる。全PEに工学用途がないという意味にしない。

\Needspace{4\baselineskip}\paragraph{選択肢の検討（査読待ち初稿）}\mbox{}\par

\noindent\textbf{1：選ぶ候補} PE。通常PEは汎用プラスチックとして選ぶ。ただし超高分子量等の特殊PEは工学用途がある。\par

\noindent\textbf{2：選ばない候補} PC。代表汎用エンプラ。\par

\noindent\textbf{3：選ばない候補} ポリアミド。ナイロン等は代表汎用エンプラ。\par

\noindent\textbf{4：選ばない候補} POM。代表汎用エンプラ。\par

\noindent\textbf{5：選ばない候補} PPO。PPEとも呼ぶ系統で代表エンプラに入る。\par

\paragraph{再利用する考え方} 物質名を暗記する前に、結合の構造、反応の種類、成立条件、数量の保存を対応させる。

\paragraph{共通理論} \hyperref[theory:organic:engineering]{第\ref*{theory:organic:engineering}節}

\par 参考：\cite{organic_enpla}。

\par 難易度 2（校準初稿）、優先度未評価。

\section{メタン四面体とエタン配座の区別}

\subsection{令和5年度 問14}
\label{q:r05-reg-applied-14}

\noindent\textbf{原典：} sources/R05/17\_applied\_science.pdf、PDF物理5頁。

\textbf{公式正答：} 5（single）。

\par\textbf{検討未確定：} 水に不溶は難溶の慣用表現で厳密ゼロではない。原典を保持し溶解度の全条件未照合を明示。

\paragraph{問題の要旨} メタンの四面体構造とエタンのねじれ・重なり配座を区別する。

\Needspace{4\baselineskip}\paragraph{確認した条件・記述（独自要約）}\mbox{}\par

\noindent\textbf{判断対象：} メタンの四面体構造とエタンのねじれ・重なり配座を区別する。\par

\noindent\textbf{原典確認：} 化学式・否定・表の全5候補を保存原典24頁画像で照合。\par

\paragraph{問われる核心} メタンの四面体構造とエタンのねじれ・重なり配座を区別する。

\paragraph{必要知識} メタンCH4は一つの炭素から四つのHへ結合する四面体である。

\paragraph{考え方・正答を選ぶ理由（査読待ち初稿）} メタンCH4は一つの炭素から四つのHへ結合する四面体である。ねじれ/重なり配座は、エタンのC-C結合両端の三つずつの結合の相対回転で区別するもので、メタンにはそのC-C軸がないので5。1atmでは沸点約−161℃で−100℃は気体。水に不溶という慣用表現は溶解度が厳密にゼロという意味にしない。

\Needspace{4\baselineskip}\paragraph{選択肢の検討（査読待ち初稿）}\mbox{}\par

\noindent\textbf{1：選ばない候補} 1atmで−100℃では液化しない。1atmの沸点は約−161℃で−100℃は気体。\par

\noindent\textbf{2：選ばない候補} 2s/2p混成で単結合説明。sp3の四面体軌道をHの1sと組合せる。\par

\noindent\textbf{3：選ばない候補} 水に不溶。通常は難溶として扱う。有限の小さい溶解度はある。\par

\noindent\textbf{4：選ばない候補} 天然ガス主成分。主要燃料成分の代表。\par

\noindent\textbf{5：選ぶ候補} ねじれ型と重なり型配座。C-C単結合の両端の相対回転で扱うエタン等の配座をメタンに当てた誤り。\par

\paragraph{再利用する考え方} 物質名を暗記する前に、結合の構造、反応の種類、成立条件、数量の保存を対応させる。

\paragraph{共通理論} \hyperref[theory:organic:methane]{第\ref*{theory:organic:methane}節}

\par 参考：\cite{organic_methane,organic_isomers}。

\par 難易度 2（校準初稿）、優先度未評価。

\section{求核性・酸性・立体対称性}

\subsection{平成26年度 問17}
\label{q:h26-reg-applied-17}

\noindent\textbf{原典：} sources/H26/17\_applied\_science.pdf、PDF物理7頁。

\textbf{公式正答：} 2（single）。

\paragraph{問題の要旨} 求核性・酸性・配向性と不斉炭素を区別する。

\Needspace{4\baselineskip}\paragraph{確認した条件・記述（独自要約）}\mbox{}\par

\noindent\textbf{判断対象：} 求核性・酸性・配向性と不斉炭素を区別する。\par

\noindent\textbf{原典確認：} 化学式・否定・表の全5候補を保存原典24頁画像で照合。\par

\paragraph{問われる核心} 求核性・酸性・配向性と不斉炭素を区別する。

\paragraph{必要知識} 求核剤は電子対を与えて結合を作り、求電子剤はそれを受け取る。

\paragraph{考え方・正答を選ぶ理由（査読待ち初稿）} 求核剤は電子対を与えて結合を作り、求電子剤はそれを受け取る。H+は求電子側なので2が不適切。OH−とCl−、孤立電子対を持つ水は条件により求核剤となる。2,4ジメチルヘキサンではC2に同じメチルが二つあるが、C4の四置換基は異なり不斉中心となる。名称に分枝があるだけで全炭素を不斉としない。

\Needspace{4\baselineskip}\paragraph{選択肢の検討（査読待ち初稿）}\mbox{}\par

\noindent\textbf{1：選ばない候補} C4以上アルカンに骨格異性。C4には直鎖と分枝の二骨格。\par

\noindent\textbf{2：選ぶ候補} OH−/Cl−/H2O/H+で求核はH+。H+は電子対を受け取る求電子剤であり逆。\par

\noindent\textbf{3：選ばない候補} 電子供与基で芳香族求電子置換促進。中間体安定化で置換を促進する代表傾向。\par

\noindent\textbf{4：選ばない候補} カルボン酸/フェノールはアルコールより酸性。共役塩基の共鳴安定化等で比較できる。\par

\noindent\textbf{5：選ばない候補} 2,4ジメチルヘキサンに不斉炭素。C4はH・メチル・右エチル・左の分枝鎖の4種を持つ。\par

\paragraph{再利用する考え方} 物質名を暗記する前に、結合の構造、反応の種類、成立条件、数量の保存を対応させる。

\paragraph{共通理論} \hyperref[theory:organic:structure]{第\ref*{theory:organic:structure}節}

\par 参考：\cite{organic_chirality,organic_directing}。

\par 難易度 2（校準初稿）、優先度未評価。

\subsection{平成29年度 問16}
\label{q:h29-reg-applied-16}

\noindent\textbf{原典：} sources/H29/17\_applied\_science.pdf、PDF物理5頁。

\textbf{公式正答：} 1（single）。

\paragraph{問題の要旨} 骨格・立体・共鳴・酸性を分けて有機分子を比較する。

\Needspace{4\baselineskip}\paragraph{確認した条件・記述（独自要約）}\mbox{}\par

\noindent\textbf{判断対象：} 骨格・立体・共鳴・酸性を分けて有機分子を比較する。\par

\noindent\textbf{原典確認：} 化学式・否定・表の全5候補を保存原典24頁画像で照合。\par

\paragraph{問われる核心} 骨格・立体・共鳴・酸性を分けて有機分子を比較する。

\paragraph{必要知識} メチルシクロヘキサンでは置換炭素に接する二つの環方向が等価で、四種置換基の不斉中心にならない。

\paragraph{考え方・正答を選ぶ理由（査読待ち初稿）} メチルシクロヘキサンでは置換炭素に接する二つの環方向が等価で、四種置換基の不斉中心にならない。鏡像異性体があるという1が不適切。環反転でaxial/equatorialが入れ替わる配座と鏡像異性は違う。1,2ジメチルシクロプロパンのcis/trans、ベンゼンの非局在π、アセチレンの酸性、過酸化物のラジカルはそれぞれ別の構造根拠で説明する。

\Needspace{4\baselineskip}\paragraph{選択肢の検討（査読待ち初稿）}\mbox{}\par

\noindent\textbf{1：選ぶ候補} メチルシクロヘキサンに鏡像異性体。メチル付炭素から環の左右へ進む結合経路が等価で、不斉中心にならない。椅子形の配座異性とは別。\par

\noindent\textbf{2：選ばない候補} ベンゼンは共鳴混成体。局在した二つの式は表現で、実分子の環へπ電子が非局在。混合物ではない。\par

\noindent\textbf{3：選ばない候補} 1,2ジメチルシクロプロパンにcis/trans。小環の同じ側と反対側の置換を回転だけでは重ねられず区別できる。\par

\noindent\textbf{4：選ばない候補} アセチレンはエチレンより酸性。sp炭素のs性50\%はsp2の33\%より高く、共役塩基を安定化する。\par

\noindent\textbf{5：選ばない候補} 過酸化ベンゾイル加熱でラジカル。O-O結合の均等開裂でラジカルを作る。\par

\paragraph{再利用する考え方} 物質名を暗記する前に、結合の構造、反応の種類、成立条件、数量の保存を対応させる。

\paragraph{共通理論} \hyperref[theory:organic:structure]{第\ref*{theory:organic:structure}節}

\par 参考：\cite{organic_chirality,organic_benzene}。

\par 難易度 2（校準初稿）、優先度未評価。

\subsection{令和2年度 問18}
\label{q:r02-reg-applied-18}

\noindent\textbf{原典：} sources/R02/17\_applied\_science.pdf、PDF物理6頁。

\textbf{公式正答：} 3（single）。

\paragraph{問題の要旨} 骨格・立体・共鳴・酸性を分けて有機分子を比較する。

\Needspace{4\baselineskip}\paragraph{確認した条件・記述（独自要約）}\mbox{}\par

\noindent\textbf{判断対象：} 骨格・立体・共鳴・酸性を分けて有機分子を比較する。\par

\noindent\textbf{原典確認：} 化学式・否定・表の全5候補を保存原典24頁画像で照合。\par

\paragraph{問われる核心} 骨格・立体・共鳴・酸性を分けて有機分子を比較する。

\paragraph{必要知識} メチルシクロヘキサンでは置換炭素に接する二つの環方向が等価で、四種置換基の不斉中心にならない。

\paragraph{考え方・正答を選ぶ理由（査読待ち初稿）} メチルシクロヘキサンでは置換炭素に接する二つの環方向が等価で、四種置換基の不斉中心にならない。鏡像異性体があるという3が不適切。環反転でaxial/equatorialが入れ替わる配座と鏡像異性は違う。1,2ジメチルシクロプロパンのcis/trans、ベンゼンの非局在π、アセチレンの酸性、過酸化物のラジカルはそれぞれ別の構造根拠で説明する。

\Needspace{4\baselineskip}\paragraph{選択肢の検討（査読待ち初稿）}\mbox{}\par

\noindent\textbf{1：選ばない候補} 過酸化ベンゾイル加熱でラジカル。O-O結合の均等開裂でラジカルを作る。\par

\noindent\textbf{2：選ばない候補} 1,2ジメチルシクロプロパンにcis/trans。小環の同じ側と反対側の置換を回転だけでは重ねられず区別できる。\par

\noindent\textbf{3：選ぶ候補} メチルシクロヘキサンに鏡像異性体。メチル付炭素から環の左右へ進む結合経路が等価で、不斉中心にならない。椅子形の配座異性とは別。\par

\noindent\textbf{4：選ばない候補} アセチレンはエチレンより酸性。sp炭素のs性50\%はsp2の33\%より高く、共役塩基を安定化する。\par

\noindent\textbf{5：選ばない候補} ベンゼンは共鳴混成体。局在した二つの式は表現で、実分子の環へπ電子が非局在。混合物ではない。\par

\paragraph{再利用する考え方} 物質名を暗記する前に、結合の構造、反応の種類、成立条件、数量の保存を対応させる。

\paragraph{共通理論} \hyperref[theory:organic:structure]{第\ref*{theory:organic:structure}節}

\par 参考：\cite{organic_chirality,organic_benzene}。

\par 難易度 2（校準初稿）、優先度未評価。

\section{石油留分・分解・改質とオクタン価}

\subsection{平成23年度 問14}
\label{q:h23-reg-applied-14}

\noindent\textbf{原典：} sources/H23/17\_applied\_science.pdf、PDF物理5頁。

\textbf{公式正答：} 4（single）。

\par\textbf{検討未確定：} 100超の出題当時測定規格全文と直鎖同族列全数値は未照合。0–100の体積百分率の説明を100超へ無条件拡張しない。

\paragraph{問題の要旨} オクタン価の標準と芳香族のアンチノック性を比較する。

\Needspace{4\baselineskip}\paragraph{確認した条件・記述（独自要約）}\mbox{}\par

\noindent\textbf{判断対象：} オクタン価の標準と芳香族のアンチノック性を比較する。\par

\noindent\textbf{原典確認：} 化学式・否定・表の全5候補を保存原典24頁画像で照合。\par

\paragraph{問われる核心} オクタン価の標準と芳香族のアンチノック性を比較する。

\paragraph{必要知識} オクタン価は燃料のアンチノック性を標準と比較する尺度で、発熱量や純度ではない。

\paragraph{考え方・正答を選ぶ理由（査読待ち初稿）} オクタン価は燃料のアンチノック性を標準と比較する尺度で、発熱量や純度ではない。nヘプタン0、イソオクタン100を置き、0–100で同等な混合標準のイソオクタン体積率で表す。芳香族への改質は高オクタン化へ用いられるので4が不適切。3は歴史的100超の記述として原典を保持し、鉛添加を現在の推奨へ変えない。

\Needspace{4\baselineskip}\paragraph{選択肢の検討（査読待ち初稿）}\mbox{}\par

\noindent\textbf{1：選ばない候補} nヘプタン/イソオクタン標準。0と100の標準物質を正しく示す。\par

\noindent\textbf{2：選ばない候補} 直鎖は炭素数が少ないほど高い。同族列の代表傾向であり混合燃料の全性質を炭素数だけで保証しない。\par

\noindent\textbf{3：選ばない候補} 100超はイソオクタンに四エチル鉛。歴史的100超の尺度の記述。出題当時規格本文は未照合。\par

\noindent\textbf{4：選ぶ候補} 芳香族は低くほぼ負。芳香族は一般に高いオクタン価へ寄与し、負が一般ではない。\par

\noindent\textbf{5：選ばない候補} 同等標準のイソオクタン体積率。0–100の標準混合物の定義である。100超へ百分率をそのまま延ばさない。\par

\paragraph{再利用する考え方} 物質名を暗記する前に、結合の構造、反応の種類、成立条件、数量の保存を対応させる。

\paragraph{共通理論} \hyperref[theory:organic:petroleum]{第\ref*{theory:organic:petroleum}節}

\par 参考：\cite{organic_gasoline}。

\par 難易度 2（校準初稿）、優先度未評価。

\subsection{平成24年度 問18}
\label{q:h24-reg-applied-18}

\noindent\textbf{原典：} sources/H24/17\_applied\_science.pdf、PDF物理6頁。

\textbf{公式正答：} 3（single）。

\par\textbf{検討未確定：} 原典の750℃を全装置の固定温度としない。水蒸気改質と酸素による部分酸化は別反応で、全工程条件・当時触媒仕様の全文照合は未完了。

\paragraph{問題の要旨} 石油留分の分解・改質・選択酸化を製品へ対応させる。

\Needspace{4\baselineskip}\paragraph{確認した条件・記述（独自要約）}\mbox{}\par

\noindent\textbf{判断対象：} 石油留分の分解・改質・選択酸化を製品へ対応させる。\par

\noindent\textbf{原典確認：} 化学式・否定・表の全5候補を保存原典24頁画像で照合。\par

\paragraph{問われる核心} 石油留分の分解・改質・選択酸化を製品へ対応させる。

\paragraph{必要知識} 分留は沸点で分け、分解・改質は結合を変える。

\paragraph{考え方・正答を選ぶ理由（査読待ち初稿）} 分留は沸点で分け、分解・改質は結合を変える。ナフサの熱分解からエチレン、プロピレン、ブタジエン等を得るので3のパラフィンが不適切。実際の分解生成物は混合物で、飽和ガスが一切ないという意味ではない。接触改質の芳香族と選択酸化・CO/H2合成ガスを別反応として整理する。

\Needspace{4\baselineskip}\paragraph{選択肢の検討（査読待ち初稿）}\mbox{}\par

\noindent\textbf{1：選ばない候補} 固体酸の接触分解で高オクタンガソリン。大きな分子の切断でガソリン原料へ。\par

\noindent\textbf{2：選ばない候補} 白金/固体酸改質で芳香族。芳香族化と異性化などによる改質。\par

\noindent\textbf{3：選ぶ候補} ナフサ熱分解でパラフィンガス。主目的のエチレン等は二重結合のあるオレフィンであり語が違う。\par

\noindent\textbf{4：選ばない候補} 完全酸化を避け選択部分酸化。目的有機化合物を残すため選択性が必要。\par

\noindent\textbf{5：選ばない候補} 合成ガスはCO/H2。メタノール等の原料となる混合ガス。\par

\paragraph{再利用する考え方} 物質名を暗記する前に、結合の構造、反応の種類、成立条件、数量の保存を対応させる。

\paragraph{共通理論} \hyperref[theory:organic:petroleum]{第\ref*{theory:organic:petroleum}節}

\par 参考：\cite{organic_jpca,organic_gasoline}。

\par 難易度 2（校準初稿）、優先度未評価。

\subsection{平成25年度 問14}
\label{q:h25-reg-applied-14}

\noindent\textbf{原典：} sources/H25/17\_applied\_science.pdf、PDF物理6頁。

\textbf{公式正答：} 1（single）。

\paragraph{問題の要旨} 石油化学原料の留分位置と反応生成物を対応させる。

\Needspace{4\baselineskip}\paragraph{確認した条件・記述（独自要約）}\mbox{}\par

\noindent\textbf{判断対象：} 石油化学原料の留分位置と反応生成物を対応させる。\par

\noindent\textbf{原典確認：} 化学式・否定・表の全5候補を保存原典24頁画像で照合。\par

\paragraph{問われる核心} 石油化学原料の留分位置と反応生成物を対応させる。

\paragraph{必要知識} ナフサは原油の低沸点側の軽質留分であり、軽油より高沸点とする1が逆である。

\paragraph{考え方・正答を選ぶ理由（査読待ち初稿）} ナフサは原油の低沸点側の軽質留分であり、軽油より高沸点とする1が逆である。分留の位置を先に固定し、エチレン→酸化→エチレンオキシド、ブタジエン→ゴム、pキシレン→テレフタル酸→PETという反応・用途の鎖を分けて覚える。

\Needspace{4\baselineskip}\paragraph{選択肢の検討（査読待ち初稿）}\mbox{}\par

\noindent\textbf{1：選ぶ候補} ナフサは軽油より高沸点。ナフサは軽質で軽油より低沸点。\par

\noindent\textbf{2：選ばない候補} ナフサ熱分解からエチレン。二重結合を持つ主要基礎製品。\par

\noindent\textbf{3：選ばない候補} エチレン酸化からエチレンオキシド。環状エーテルの酸化生成物。\par

\noindent\textbf{4：選ばない候補} 1,3ブタジエンは合成ゴム原料。ジエンの重合をゴムへ利用。\par

\noindent\textbf{5：選ばない候補} pキシレン酸化からPET原料酸。テレフタル酸をエチレングリコールとポリエステルへ。\par

\paragraph{再利用する考え方} 物質名を暗記する前に、結合の構造、反応の種類、成立条件、数量の保存を対応させる。

\paragraph{共通理論} \hyperref[theory:organic:petroleum]{第\ref*{theory:organic:petroleum}節}

\par 参考：\cite{organic_jpca_dict,organic_step}。

\par 難易度 2（校準初稿）、優先度未評価。

\section{ナイロン6/66の結合と製造経路}

\subsection{平成23年度 問16}
\label{q:h23-reg-applied-16}

\noindent\textbf{原典：} sources/H23/17\_applied\_science.pdf、PDF物理6頁。

\textbf{公式正答：} 3（single）。

\paragraph{問題の要旨} ナイロン6と66の原料経路を区別する。

\Needspace{4\baselineskip}\paragraph{確認した条件・記述（独自要約）}\mbox{}\par

\noindent\textbf{判断対象：} ナイロン6と66の原料経路を区別する。\par

\noindent\textbf{原典確認：} 化学式・否定・表の全5候補を保存原典24頁画像で照合。\par

\paragraph{問われる核心} ナイロン6と66の原料経路を区別する。

\paragraph{必要知識} ナイロン6はカプロラクタムの開環を利用するポリアミドである。

\paragraph{考え方・正答を選ぶ理由（査読待ち初稿）} ナイロン6はカプロラクタムの開環を利用するポリアミドである。代表経路はシクロヘキサン→シクロヘキサノン→オキシム→カプロラクタム。アジピン酸はヘキサメチレンジアミンとナイロン66を作る原料なので3を選ぶ。ポリアミドという結合種が共通でも原料は共通でない。

\Needspace{4\baselineskip}\paragraph{選択肢の検討（査読待ち初稿）}\mbox{}\par

\noindent\textbf{1：選ばない候補} シクロヘキサン。環状六炭素の出発原料経路に属す。\par

\noindent\textbf{2：選ばない候補} シクロヘキサノン。オキシムの前駆体。\par

\noindent\textbf{3：選ぶ候補} アジピン酸。ジアミンとの縮合でナイロン66を作る原料。\par

\noindent\textbf{4：選ばない候補} シクロヘキサノンオキシム。Beckmann転位でラクタムへ進む。\par

\noindent\textbf{5：選ばない候補} εカプロラクタム。環を開いてナイロン6のポリアミド鎖になる。\par

\paragraph{再利用する考え方} 物質名を暗記する前に、結合の構造、反応の種類、成立条件、数量の保存を対応させる。

\paragraph{共通理論} \hyperref[theory:organic:nylon]{第\ref*{theory:organic:nylon}節}

\par 参考：\cite{organic_step}。

\par 難易度 2（校準初稿）、優先度未評価。

\section{高分子の架橋と熱応答}

\subsection{平成23年度 問9}
\label{q:h23-reg-applied-09}

\noindent\textbf{原典：} sources/H23/17\_applied\_science.pdf、PDF物理4頁。

\textbf{公式正答：} 1（single）。

\paragraph{問題の要旨} 架橋による熱硬化と鎖の流動による熱可塑を区別する。

\Needspace{4\baselineskip}\paragraph{確認した条件・記述（独自要約）}\mbox{}\par

\noindent\textbf{判断対象：} 架橋による熱硬化と鎖の流動による熱可塑を区別する。\par

\noindent\textbf{原典確認：} 化学式・否定・表の全5候補を保存原典24頁画像で照合。\par

\paragraph{問われる核心} 架橋による熱硬化と鎖の流動による熱可塑を区別する。

\paragraph{必要知識} 熱硬化は加熱で鎖が流れることではなく、反応で架橋網が形成されることである。

\paragraph{考え方・正答を選ぶ理由（査読待ち初稿）} 熱硬化は加熱で鎖が流れることではなく、反応で架橋網が形成されることである。フェノールとホルムアルデヒドの樹脂は架橋した熱硬化性で1。残る通常のPS・PVC・PMMA・PVDCは熱可塑性である。熱硬化性でも高温で化学分解するので、無限耐熱ではない。

\Needspace{4\baselineskip}\paragraph{選択肢の検討（査読待ち初稿）}\mbox{}\par

\noindent\textbf{1：選ぶ候補} フェノール・ホルムアルデヒド。三次元架橋網を作り硬化後は再溶融しない。\par

\noindent\textbf{2：選ばない候補} ポリスチレン。通常のPSは熱可塑。\par

\noindent\textbf{3：選ばない候補} PVC。通常のPVCは熱可塑。\par

\noindent\textbf{4：選ばない候補} PMMA。通常のPMMAは熱可塑。\par

\noindent\textbf{5：選ばない候補} PVDC。通常のPVDCは熱可塑。\par

\paragraph{再利用する考え方} 物質名を暗記する前に、結合の構造、反応の種類、成立条件、数量の保存を対応させる。

\paragraph{共通理論} \hyperref[theory:organic:polymer]{第\ref*{theory:organic:polymer}節}

\par 参考：\cite{organic_physical,organic_enpla}。

\par 難易度 2（校準初稿）、優先度未評価。

\subsection{平成28年度 問17}
\label{q:h28-reg-applied-17}

\noindent\textbf{原典：} sources/H28/17\_applied\_science.pdf、PDF物理6頁。

\textbf{公式正答：} 4（single）。

\par\textbf{検討未確定：} ナイロン全般をジカルボン酸/ジアミン由来へ限定する候補はナイロン6等に当てはまらない。公式PVC肢の説明と別に原典の一般化疑義を残す。

\paragraph{問題の要旨} 高分子の結合・熱応答・生分解性を別軸で比較する。

\Needspace{4\baselineskip}\paragraph{確認した条件・記述（独自要約）}\mbox{}\par

\noindent\textbf{判断対象：} 高分子の結合・熱応答・生分解性を別軸で比較する。\par

\noindent\textbf{原典確認：} 化学式・否定・表の全5候補を保存原典24頁画像で照合。\par

\paragraph{問われる核心} 高分子の結合・熱応答・生分解性を別軸で比較する。

\paragraph{必要知識} PVCは通常の熱可塑性樹脂であり、熱硬化性という候補4が公式の不適切肢である。

\paragraph{考え方・正答を選ぶ理由（査読待ち初稿）} PVCは通常の熱可塑性樹脂であり、熱硬化性という候補4が公式の不適切肢である。タンパク質・ナイロンはアミド、PET・PLAはエステルを主鎖に持つが、結合種だけで熱応答と生分解性を同一視しない。原典のナイロンの二原料縮合はナイロン66等へ限定すべきで、ナイロン6を含む全称では正確でない。この留保も保持する。

\Needspace{4\baselineskip}\paragraph{選択肢の検討（査読待ち初稿）}\mbox{}\par

\noindent\textbf{1：選ばない候補} タンパク質はアミノ酸の生体高分子。アミノ酸残基がアミド結合でつながる生体高分子。\par

\noindent\textbf{2：選ばない候補} ナイロンはジカルボン酸/ジアミンのポリアミド。ナイロン66なら二官能原料の縮合が対応するが、全ナイロンの定義ではなくナイロン6はラクタム開環。\par

\noindent\textbf{3：選ばない候補} PLAは生分解性。代表的生分解性ポリエステル。ただし環境条件を伴う。\par

\noindent\textbf{4：選ぶ候補} PVCは熱硬化性。通常PVCは熱可塑性で熱硬化性ではない。\par

\noindent\textbf{5：選ばない候補} PETは熱可塑性。通常PETは熱可塑性ポリエステル。\par

\paragraph{再利用する考え方} 物質名を暗記する前に、結合の構造、反応の種類、成立条件、数量の保存を対応させる。

\paragraph{共通理論} \hyperref[theory:organic:polymer]{第\ref*{theory:organic:polymer}節}

\par 参考：\cite{organic_physical,organic_step,organic_jbpa}。

\par 難易度 2（校準初稿）、優先度未評価。

\subsection{令和6年度 問17}
\label{q:r06-reg-applied-17}

\noindent\textbf{原典：} sources/R06/17\_applied\_science.pdf、PDF物理5頁。

\textbf{公式正答：} 3（single）。

\par\textbf{検討未確定：} ナイロン全般をジカルボン酸/ジアミン由来へ限定する候補はナイロン6等に当てはまらない。公式PVC肢の説明と別に原典の一般化疑義を残す。

\paragraph{問題の要旨} 高分子の結合・熱応答・生分解性を別軸で比較する。

\Needspace{4\baselineskip}\paragraph{確認した条件・記述（独自要約）}\mbox{}\par

\noindent\textbf{判断対象：} 高分子の結合・熱応答・生分解性を別軸で比較する。\par

\noindent\textbf{原典確認：} 化学式・否定・表の全5候補を保存原典24頁画像で照合。\par

\paragraph{問われる核心} 高分子の結合・熱応答・生分解性を別軸で比較する。

\paragraph{必要知識} PVCは通常の熱可塑性樹脂であり、熱硬化性という候補3が公式の不適切肢である。

\paragraph{考え方・正答を選ぶ理由（査読待ち初稿）} PVCは通常の熱可塑性樹脂であり、熱硬化性という候補3が公式の不適切肢である。タンパク質・ナイロンはアミド、PET・PLAはエステルを主鎖に持つが、結合種だけで熱応答と生分解性を同一視しない。原典のナイロンの二原料縮合はナイロン66等へ限定すべきで、ナイロン6を含む全称では正確でない。この留保も保持する。

\Needspace{4\baselineskip}\paragraph{選択肢の検討（査読待ち初稿）}\mbox{}\par

\noindent\textbf{1：選ばない候補} ナイロンはジカルボン酸/ジアミンのポリアミド。ナイロン66なら二官能原料の縮合が対応するが、全ナイロンの定義ではなくナイロン6はラクタム開環。\par

\noindent\textbf{2：選ばない候補} PLAは生分解性。代表的生分解性ポリエステル。ただし環境条件を伴う。\par

\noindent\textbf{3：選ぶ候補} PVCは熱硬化性。通常PVCは熱可塑性で熱硬化性ではない。\par

\noindent\textbf{4：選ばない候補} PETは熱可塑性。通常PETは熱可塑性ポリエステル。\par

\noindent\textbf{5：選ばない候補} タンパク質はアミノ酸の生体高分子。アミノ酸残基がアミド結合でつながる生体高分子。\par

\paragraph{再利用する考え方} 物質名を暗記する前に、結合の構造、反応の種類、成立条件、数量の保存を対応させる。

\paragraph{共通理論} \hyperref[theory:organic:polymer]{第\ref*{theory:organic:polymer}節}

\par 参考：\cite{organic_physical,organic_step,organic_jbpa}。

\par 難易度 2（校準初稿）、優先度未評価。

\section{有機物の密度・揮発性と反応性}

\subsection{平成27年度 問14}
\label{q:h27-reg-applied-14}

\noindent\textbf{原典：} sources/H27/17\_applied\_science.pdf、PDF物理5頁。

\textbf{公式正答：} 4（single）。

\par\textbf{検討未確定：} 過酸化物の加熱・摩擦を全条件で必ず爆発するとの全称にしない。各候補の出題時SDS・全条件は未照合、公開ICSCの対象性状を参照。

\paragraph{問題の要旨} 物質ごとの密度・揮発性・反応性を比較する。

\Needspace{4\baselineskip}\paragraph{確認した条件・記述（独自要約）}\mbox{}\par

\noindent\textbf{判断対象：} 物質ごとの密度・揮発性・反応性を比較する。\par

\noindent\textbf{原典確認：} 化学式・否定・表の全5候補を保存原典24頁画像で照合。\par

\paragraph{問われる核心} 物質ごとの密度・揮発性・反応性を比較する。

\paragraph{必要知識} ニトロベンゼンの相対密度は水=1に対して約1.2である。

\paragraph{考え方・正答を選ぶ理由（査読待ち初稿）} ニトロベンゼンの相対密度は水=1に対して約1.2である。水より小さいとする候補4が不適切。エーテルの揮発・引火、ベンゼン蒸気の有害性、過酸化物の分解、アクリル酸の重合は異なる物性・危険である。爆発や重合は状態・温度・組成に依存し、原典の一般表現を全条件での必然としない。

\Needspace{4\baselineskip}\paragraph{選択肢の検討（査読待ち初稿）}\mbox{}\par

\noindent\textbf{1：選ばない候補} ジエチルエーテル揮発・引火。低沸点有機溶媒で蒸気が燃えやすい。\par

\noindent\textbf{2：選ばない候補} 過酸化ベンゾイル加熱・摩擦で爆発分解。弱いO-O結合の分解と発熱・感度を持ち、爆発的分解の危険がある。必ず爆発するという無条件則ではない。\par

\noindent\textbf{3：選ばない候補} ベンゼン揮発・有毒蒸気。蒸発して有害な暴露を起こし得る。\par

\noindent\textbf{4：選ぶ候補} ニトロベンゼンは水より軽い。公開ICSCの水=1に対する比重は約1.2で、軽いという記述が逆。\par

\noindent\textbf{5：選ばない候補} アクリル酸は重合しやすい。ビニル結合の重合で発熱し、条件・阻害剤の管理が必要。\par

\paragraph{再利用する考え方} 物質名を暗記する前に、結合の構造、反応の種類、成立条件、数量の保存を対応させる。

\paragraph{共通理論} \hyperref[theory:organic:properties]{第\ref*{theory:organic:properties}節}

\par 参考：\cite{organic_nitro,organic_peroxide}。

\par 難易度 2（校準初稿）、優先度未評価。

\subsection{令和3年度 問12}
\label{q:r03-reg-applied-12}

\noindent\textbf{原典：} sources/R03/17\_applied\_science.pdf、PDF物理5頁。

\textbf{公式正答：} 2（single）。

\par\textbf{検討未確定：} 過酸化物の加熱・摩擦を全条件で必ず爆発するとの全称にしない。各候補の出題時SDS・全条件は未照合、公開ICSCの対象性状を参照。

\paragraph{問題の要旨} 物質ごとの密度・揮発性・反応性を比較する。

\Needspace{4\baselineskip}\paragraph{確認した条件・記述（独自要約）}\mbox{}\par

\noindent\textbf{判断対象：} 物質ごとの密度・揮発性・反応性を比較する。\par

\noindent\textbf{原典確認：} 化学式・否定・表の全5候補を保存原典24頁画像で照合。\par

\paragraph{問われる核心} 物質ごとの密度・揮発性・反応性を比較する。

\paragraph{必要知識} ニトロベンゼンの相対密度は水=1に対して約1.2である。

\paragraph{考え方・正答を選ぶ理由（査読待ち初稿）} ニトロベンゼンの相対密度は水=1に対して約1.2である。水より小さいとする候補2が不適切。エーテルの揮発・引火、ベンゼン蒸気の有害性、過酸化物の分解、アクリル酸の重合は異なる物性・危険である。爆発や重合は状態・温度・組成に依存し、原典の一般表現を全条件での必然としない。

\Needspace{4\baselineskip}\paragraph{選択肢の検討（査読待ち初稿）}\mbox{}\par

\noindent\textbf{1：選ばない候補} ベンゼン揮発・有毒蒸気。蒸発して有害な暴露を起こし得る。\par

\noindent\textbf{2：選ぶ候補} ニトロベンゼンは水より軽い。公開ICSCの水=1に対する比重は約1.2で、軽いという記述が逆。\par

\noindent\textbf{3：選ばない候補} 過酸化ベンゾイル加熱・摩擦で爆発分解。弱いO-O結合の分解と発熱・感度を持ち、爆発的分解の危険がある。必ず爆発するという無条件則ではない。\par

\noindent\textbf{4：選ばない候補} アクリル酸は重合しやすい。ビニル結合の重合で発熱し、条件・阻害剤の管理が必要。\par

\noindent\textbf{5：選ばない候補} ジエチルエーテル揮発・引火。低沸点有機溶媒で蒸気が燃えやすい。\par

\paragraph{再利用する考え方} 物質名を暗記する前に、結合の構造、反応の種類、成立条件、数量の保存を対応させる。

\paragraph{共通理論} \hyperref[theory:organic:properties]{第\ref*{theory:organic:properties}節}

\par 参考：\cite{organic_nitro,organic_peroxide}。

\par 難易度 2（校準初稿）、優先度未評価。

\section{両親媒性・界面吸着・ミセルとCMC}

\subsection{平成27年度 問13}
\label{q:h27-reg-applied-13}

\noindent\textbf{原典：} sources/H27/17\_applied\_science.pdf、PDF物理5頁。

\textbf{公式正答：} 2（single）。

\paragraph{問題の要旨} 界面活性分子と固体吸着材を区別する。

\Needspace{4\baselineskip}\paragraph{確認した条件・記述（独自要約）}\mbox{}\par

\noindent\textbf{判断対象：} 界面活性分子と固体吸着材を区別する。\par

\noindent\textbf{原典確認：} 化学式・否定・表の全5候補を保存原典24頁画像で照合。\par

\paragraph{問われる核心} 界面活性分子と固体吸着材を区別する。

\paragraph{必要知識} 界面活性剤は親水部と疎水部を同一分子に持つ。

\paragraph{考え方・正答を選ぶ理由（査読待ち初稿）} 界面活性剤は親水部と疎水部を同一分子に持つ。けん化は石けん生成、CMCは会合、ソルビタン脂肪酸エステルとアルキルベンゼンスルホン酸塩は代表剤に対応する。活性炭は多孔質固体への吸着に用いられるため、この五候補で最も関係が薄い2。活性炭に界面がないとか吸着しないという説明ではない。

\Needspace{4\baselineskip}\paragraph{選択肢の検討（査読待ち初稿）}\mbox{}\par

\noindent\textbf{1：選ばない候補} けん化。脂肪酸塩の石けんを生成。\par

\noindent\textbf{2：選ぶ候補} 活性炭。多孔質の固体吸着材で界面活性分子とは別分類。\par

\noindent\textbf{3：選ばない候補} 臨界ミセル濃度。界面活性分子の会合が重要になる濃度。\par

\noindent\textbf{4：選ばない候補} ソルビタン脂肪酸エステル。非イオン系の代表。\par

\noindent\textbf{5：選ばない候補} ドデシルベンゼンスルホン酸Na。陰イオン系の代表。\par

\paragraph{再利用する考え方} 物質名を暗記する前に、結合の構造、反応の種類、成立条件、数量の保存を対応させる。

\paragraph{共通理論} \hyperref[theory:organic:surfactant]{第\ref*{theory:organic:surfactant}節}

\par 参考：\cite{organic_cmc}。

\par 難易度 2（校準初稿）、優先度未評価。

\subsection{平成29年度 問17}
\label{q:h29-reg-applied-17}

\noindent\textbf{原典：} sources/H29/17\_applied\_science.pdf、PDF物理5頁。

\textbf{公式正答：} 2（single）。

\paragraph{問題の要旨} 界面活性分子と固体吸着材を区別する。

\Needspace{4\baselineskip}\paragraph{確認した条件・記述（独自要約）}\mbox{}\par

\noindent\textbf{判断対象：} 界面活性分子と固体吸着材を区別する。\par

\noindent\textbf{原典確認：} 化学式・否定・表の全5候補を保存原典24頁画像で照合。\par

\paragraph{問われる核心} 界面活性分子と固体吸着材を区別する。

\paragraph{必要知識} 界面活性剤は親水部と疎水部を同一分子に持つ。

\paragraph{考え方・正答を選ぶ理由（査読待ち初稿）} 界面活性剤は親水部と疎水部を同一分子に持つ。けん化は石けん生成、CMCは会合、ソルビタン脂肪酸エステルとアルキルベンゼンスルホン酸塩は代表剤に対応する。活性炭は多孔質固体への吸着に用いられるため、この五候補で最も関係が薄い2。活性炭に界面がないとか吸着しないという説明ではない。

\Needspace{4\baselineskip}\paragraph{選択肢の検討（査読待ち初稿）}\mbox{}\par

\noindent\textbf{1：選ばない候補} けん化。脂肪酸塩の石けんを生成。\par

\noindent\textbf{2：選ぶ候補} 活性炭。多孔質の固体吸着材で界面活性分子とは別分類。\par

\noindent\textbf{3：選ばない候補} 臨界ミセル濃度。界面活性分子の会合が重要になる濃度。\par

\noindent\textbf{4：選ばない候補} ソルビタン脂肪酸エステル。非イオン系の代表。\par

\noindent\textbf{5：選ばない候補} ドデシルベンゼンスルホン酸Na。陰イオン系の代表。\par

\paragraph{再利用する考え方} 物質名を暗記する前に、結合の構造、反応の種類、成立条件、数量の保存を対応させる。

\paragraph{共通理論} \hyperref[theory:organic:surfactant]{第\ref*{theory:organic:surfactant}節}

\par 参考：\cite{organic_cmc}。

\par 難易度 2（校準初稿）、優先度未評価。

\subsection{令和2年度 問15}
\label{q:r02-reg-applied-15}

\noindent\textbf{原典：} sources/R02/17\_applied\_science.pdf、PDF物理5頁。

\textbf{公式正答：} 2（single）。

\par\textbf{検討未確定：} 原典の脂質全般という表現はリン脂質などへ限定して説明。全脂質の強い界面活性を保証しない。

\paragraph{問題の要旨} 界面吸着と溶液中の単量体・会合体を区別する。

\Needspace{4\baselineskip}\paragraph{確認した条件・記述（独自要約）}\mbox{}\par

\noindent\textbf{判断対象：} 界面吸着と溶液中の単量体・会合体を区別する。\par

\noindent\textbf{原典確認：} 化学式・否定・表の全5候補を保存原典24頁画像で照合。\par

\paragraph{問われる核心} 界面吸着と溶液中の単量体・会合体を区別する。

\paragraph{必要知識} 界面への吸着と溶液中の存在は両立する。

\paragraph{考え方・正答を選ぶ理由（査読待ち初稿）} 界面への吸着と溶液中の存在は両立する。CMC付近以下にも単量体があり、以上では追加分が主に会合体へ入り得る。界面以外の濃度を無視するとする2は不適切。脂質の代表はリン脂質等の両親媒性分子であり、全ての脂質を同じ界面活性剤とする説明は避ける。

\Needspace{4\baselineskip}\paragraph{選択肢の検討（査読待ち初稿）}\mbox{}\par

\noindent\textbf{1：選ばない候補} 親水/疎水部分を持つ。両親媒性の構造。\par

\noindent\textbf{2：選ぶ候補} 界面以外の濃度は無視できる。吸着選択性があってもバルクの単量体やミセル濃度は無視できない。\par

\noindent\textbf{3：選ばない候補} 脂質は天然界面活性物質。リン脂質等の両親媒性脂質に当てはまる代表的説明。全脂質一律ではない。\par

\noindent\textbf{4：選ばない候補} 水油の界面エネルギー低下。界面へ配向して張力を下げる。\par

\noindent\textbf{5：選ばない候補} 洗剤以外にも乳化・消毒・帯電防止。構造や電荷に応じた利用がある。全剤が全用途を持つとはしない。\par

\paragraph{再利用する考え方} 物質名を暗記する前に、結合の構造、反応の種類、成立条件、数量の保存を対応させる。

\paragraph{共通理論} \hyperref[theory:organic:surfactant]{第\ref*{theory:organic:surfactant}節}

\par 参考：\cite{organic_cmc}。

\par 難易度 2（校準初稿）、優先度未評価。

\subsection{令和4年度 問16}
\label{q:r04-reg-applied-16}

\noindent\textbf{原典：} sources/R04/17\_applied\_science.pdf、PDF物理5頁。

\textbf{公式正答：} 4（single）。

\par\textbf{検討未確定：} 脂質全般を界面活性剤とする原典の全称に留保。CMCと単量体の説明は溶媒・温度・濃度の条件を要する。

\paragraph{問題の要旨} ミセル形成の濃度条件から単量体の存在を判定する。

\Needspace{4\baselineskip}\paragraph{確認した条件・記述（独自要約）}\mbox{}\par

\noindent\textbf{判断対象：} ミセル形成の濃度条件から単量体の存在を判定する。\par

\noindent\textbf{原典確認：} 化学式・否定・表の全5候補を保存原典24頁画像で照合。\par

\paragraph{問われる核心} ミセル形成の濃度条件から単量体の存在を判定する。

\paragraph{必要知識} 会合は濃度等の条件に依存する。

\paragraph{考え方・正答を選ぶ理由（査読待ち初稿）} 会合は濃度等の条件に依存する。CMCより低い濃度で単量体を持ち、ミセル形成後も単量体がゼロになるわけではない。常に会合し単量体なしという4が不適切。水油界面への吸着とバルクの存在は別の平衡である。原典の脂質を代表的両親媒性脂質へ限定する留保も追跡する。

\Needspace{4\baselineskip}\paragraph{選択肢の検討（査読待ち初稿）}\mbox{}\par

\noindent\textbf{1：選ばない候補} 水油界面エネルギー低下。界面吸着による典型作用。\par

\noindent\textbf{2：選ばない候補} 脂質は天然界面活性剤。リン脂質等の両親媒性脂質を代表とする。全脂質全称に留保。\par

\noindent\textbf{3：選ばない候補} 乳化・消毒・帯電防止にも利用。目的に適した種類を選ぶ利用。\par

\noindent\textbf{4：選ぶ候補} 常に会合し単量体は存在しない。CMC未満で単量体が存在するので常に/ないが誤り。\par

\noindent\textbf{5：選ばない候補} 親水/疎水部を持つ。両親媒性の基本構造。\par

\paragraph{再利用する考え方} 物質名を暗記する前に、結合の構造、反応の種類、成立条件、数量の保存を対応させる。

\paragraph{共通理論} \hyperref[theory:organic:surfactant]{第\ref*{theory:organic:surfactant}節}

\par 参考：\cite{organic_cmc}。

\par 難易度 2（校準初稿）、優先度未評価。
